\documentclass[11pt]{amsart}

\usepackage[T1]{fontenc}
\usepackage{lmodern}
\usepackage{microtype}
\usepackage{amsmath,amssymb,amsthm}
\usepackage[hidelinks]{hyperref}

\hypersetup{
  pdftitle={A Counterexample to Kourovka Problem 20.28}
}

\newtheorem{theorem}{Theorem}
\newtheorem{proposition}[theorem]{Proposition}
\newtheorem{corollary}[theorem]{Corollary}
\theoremstyle{remark}
\newtheorem{remark}[theorem]{Remark}

\newcommand{\cs}{\operatorname{cs}}
\newcommand{\F}{\mathbb F}

\title{A Counterexample to Kourovka Problem 20.28}
\subjclass[2020]{20E45, 20J06}
\keywords{conjugacy class sizes, Schur multiplier, universal central extension,
central product, isoclinism}
\date{}

\begin{document}

\begin{abstract}
We give a negative answer to Kourovka Problem~20.28.  The group
\[
G=\{A\in \mathrm{GL}_2(5):\det A\in(\F_5^\times)^2\}
\]
has the same set of conjugacy-class sizes as \(\mathrm{SL}_2(5)\), namely
\(\{1,12,20,30\}\), but is not \(\mathrm{SL}_2(5)\times A\) for any
abelian group \(A\).  Thus the problem fails already for \(L=A_5\), and a
published characterization of \(\mathrm{SL}_2(5)\) is false as stated.
More generally, the same central enlargement gives a counterexample for every
nonabelian finite simple group with nontrivial Schur multiplier.
\end{abstract}

\maketitle

For a finite group \(X\), write
\[
\cs(X)=\{|x^X|:x\in X\}.
\]
Kourovka Problem~20.28 asks whether, for a nonabelian finite simple group
\(L\) with universal perfect central extension \(H(L)\), the equality
\(\cs(G)=\cs(H(L))\) forces \(G\cong H(L)\times A\) for some abelian group
\(A\) \cite[Problem~20.28]{Kourovka}.  Gorshkov asserted this for
\(L=A_5\), where \(H(L)\cong \mathrm{SL}_2(5)\)
\cite[Theorem~1]{Gorshkov}.

\begin{theorem}\label{thm:main}
Let
\[
G=\{A\in \mathrm{GL}_2(5):\det A\in\{1,4\}\}.
\]
Then
\[
\cs(G)=\cs(\mathrm{SL}_2(5))=\{1,12,20,30\},
\]
but \(G\not\cong \mathrm{SL}_2(5)\times A\) for any abelian group
\(A\).  Consequently, Kourovka Problem~20.28 has a negative answer, already
for \(L=A_5\), and \cite[Theorem~1]{Gorshkov} is false as stated.
\end{theorem}

The construction is a central enlargement within an isoclinism family.
We record it in a form that applies to every nontrivial Schur multiplier.

\begin{proposition}\label{prop:central}
Let \(H\) be a finite perfect group with \(Z(H)\ne 1\).  Choose a prime
\(p\mid |Z(H)|\) and an element \(z\in Z(H)\) whose order \(p^a\) is
maximal among the orders of the \(p\)-elements of \(Z(H)\).  Let
\(C=\langle c\rangle\cong C_{p^{a+1}}\), and put
\[
E=\frac{H\times C}{\langle(z,c^{-p})\rangle}.
\]
Then
\[
\cs(E)=\cs(H),\qquad |E|=p|H|,\qquad E'=H,
\]
and \(E\not\cong H\times A\) for any abelian group \(A\).
\end{proposition}

\begin{proof}
The natural map \(H\to E\) is injective, so we identify \(H\) with its
image.  Let \(t\) be the image of \((1,c)\).  Then
\[
t^p=z,\qquad |t|=p^{a+1},\qquad
E=H\langle t\rangle,\qquad H\cap\langle t\rangle=\langle z\rangle.
\]
In particular, \(|E|=p|H|\).  Since \(t\) is central, for
\(h\in H\) and \(j\in\mathbb Z\),
\[
C_E(ht^j)=C_H(h)\langle t\rangle.
\]
The intersection of the two factors on the right is \(\langle z\rangle\),
so
\[
|C_E(ht^j)|=p|C_H(h)|
\quad\text{and hence}\quad
|(ht^j)^E|=|h^H|.
\]
Thus \(\cs(E)=\cs(H)\).

An element \(ht^j\) is central in \(E\) exactly when \(h\in Z(H)\), and
therefore
\[
Z(E)=Z(H)\langle t\rangle.
\]
Also \(E/H\) is cyclic, so \(E'\le H\); since \(H\) is perfect,
\(H=H'\le E'\), whence \(E'=H\).

If \(E\cong H\times A\) with \(A\) abelian, then the order equality forces
\(A\cong C_p\).  By the choice of \(z\), the largest order of a
\(p\)-element of \(Z(H)\times C_p\) is \(p^a\), whereas
\(t\in Z(E)\) has order \(p^{a+1}\).  The centers are therefore not
isomorphic, a contradiction.
\end{proof}

\begin{remark}
The groups \(H\) and \(E\) in Proposition~\ref{prop:central} are isoclinic
in the sense of Hall \cite{Hall}: one has
\(E/Z(E)\cong H/Z(H)\), \(E'=H\), and the commutator maps agree.  The
calculation above makes the equality of class-size sets explicit.  It does
not preserve multiplicities: every conjugacy class of \(H\) gives \(p\)
conjugacy classes of \(E\) of the same size.
\end{remark}

\begin{proof}[Proof of Theorem~\ref{thm:main}]
Set \(H=\mathrm{SL}_2(5)\), \(z=-I\), and \(s=2I\).  Then
\[
s^2=z,\qquad |s|=4,\qquad \det s=4.
\]
Since \((\F_5^\times)^2=\{1,4\}\), every element of \(G\) with determinant
\(4\) lies in \(Hs\).  Hence
\[
G=H\langle s\rangle,
\qquad H\cap\langle s\rangle=\langle z\rangle,
\]
and therefore
\[
G\cong
\frac{H\times C_4}{\langle(z,c^2)\rangle}.
\]
Proposition~\ref{prop:central}, with \(p=2\) and \(a=1\), gives
\(\cs(G)=\cs(H)\) and excludes an isomorphism \(G\cong H\times A\).

It remains only to record \(\cs(H)\).  We have \(|H|=120\).  For a
noncentral \(h\in H\), the characteristic polynomial is of exactly one of
the following types.  If it has a repeated root \(\varepsilon=\pm1\), then
\(h=\varepsilon I+N\) with \(N^2=0\ne N\); its centralizer algebra is
\(\{aI+bN:a,b\in\F_5\}\), and the determinant-one units number
\(2\cdot5=10\).  If the polynomial splits over \(\F_5\) with distinct
roots, then \(C_H(h)\) is a split torus of order \(4\).  If it is
irreducible over \(\F_5\), then \(\F_5[h]\cong\F_{25}\), and
\(C_H(h)\) is the kernel of the norm
\(\F_{25}^\times\to\F_5^\times\), so \(|C_H(h)|=6\).  All three cases
occur, for example for
\[
\begin{pmatrix}1&1\\0&1\end{pmatrix},\qquad
\begin{pmatrix}2&0\\0&3\end{pmatrix},\qquad
\begin{pmatrix}0&-1\\1&1\end{pmatrix},
\]
respectively.  Thus the noncentral class sizes are
\(120/10=12\), \(120/6=20\), and \(120/4=30\), proving
\[
\cs(H)=\{1,12,20,30\}.
\]

For later comparison, Proposition~\ref{prop:central} also gives
\[
G'=H,\qquad Z(G)=\langle s\rangle\cong C_4.
\]
Every nontrivial subgroup of \(Z(G)\) contains
\(s^2=z\in G'\).  Hence \(G\) has no nontrivial abelian direct factor.
\end{proof}

\begin{corollary}\label{cor:general}
Let \(L\) be a nonabelian finite simple group with nontrivial Schur
multiplier \(M(L)\), and let \(H(L)\) be its universal perfect central
extension.  Then there is a finite group \(G_L\) such that
\[
\cs(G_L)=\cs(H(L)),
\qquad
G_L\not\cong H(L)\times A
\]
for any abelian group \(A\).
\end{corollary}

\begin{proof}
The group \(H(L)\) is perfect and
\(Z(H(L))=M(L)\ne1\).  Apply Proposition~\ref{prop:central}.
\end{proof}

\section*{The gap in the published proof}

At the end of the proof of \cite[Theorem~1]{Gorshkov}, a normal subgroup
\(R\le Z(T)\) with \(T/R\cong A_5\) is obtained.  The proof then uses the
absence of abelian direct factors to identify \(R\) as a Schur-multiplier
kernel.  A covering kernel must, however, satisfy the stem condition
\[
R\le T'
\]
in addition to centrality.  The example above realizes the missing case:
with \(T=G\) and \(R=Z(G)\),
\[
T/R\cong A_5,\qquad R\cong C_4,\qquad
T'=\mathrm{SL}_2(5),\qquad R\nleq T',
\]
while \(T\) has no nontrivial abelian direct factor.  Thus the final
inference does not apply.

\begin{thebibliography}{9}

\bibitem{Gorshkov}
I.~Gorshkov,
\emph{On characterization of a finite group by the set of conjugacy class
sizes},
J. Algebra Appl. \textbf{21} (2022), no.~11, Article ID 2250226, 8 pp.,
\href{https://doi.org/10.1142/S0219498822502267}
{doi:10.1142/S0219498822502267}.

\bibitem{Hall}
P.~Hall,
\emph{The classification of prime-power groups},
J. Reine Angew. Math. \textbf{182} (1940), 130--141,
\href{https://doi.org/10.1515/crll.1940.182.130}
{doi:10.1515/crll.1940.182.130}.

\bibitem{Kourovka}
E.~I. Khukhro and V.~D. Mazurov, eds.,
\emph{Unsolved Problems in Group Theory. The Kourovka Notebook},
no.~21, Sobolev Institute of Mathematics, Novosibirsk, 2026,
\href{https://arxiv.org/abs/1401.0300v45}{arXiv:1401.0300v45}.

\end{thebibliography}

\end{document}